
% !TEX TS-program = xelatex
\documentclass[11pt]{article}

\usepackage[a4paper,margin=25mm]{geometry}
\usepackage{fontspec}
\setmainfont{TeX Gyre Heros}
\renewcommand{\familydefault}{\sfdefault}

\usepackage{amsmath,amssymb,amsthm}
\usepackage{enumitem}
\usepackage{hyperref}

\setlist[itemize]{leftmargin=*}
\setlist[enumerate]{leftmargin=*}

\title{ENEL445 --- Structured Problem-Solving Notes \\ Lecture 1 and Week 1 Lecture 2}
\author{}
\date{}

\begin{document}
\maketitle

\section*{How to Use These Notes}

These notes are ordered to follow:

\begin{itemize}
\item Lecture 1 (Foundations of Optimisation)
\item Week 1 Lecture 2 (Duality and Structure)
\end{itemize}

Each section includes a reference tag so you can map directly back to the PDF:

\begin{itemize}
\item [L1-x] refers to Lecture 1 section x
\item [W1L2-x] refers to Week 1 Lecture 2 section x
\end{itemize}

The focus is strictly on what you must know to solve problems like those shown in the sample questions.

\newpage

% ============================================================
\section*{[L1-1] Optimisation Problem Structure}

\subsection*{Standard Form}

A general constrained optimisation problem:

\begin{align*}
\min_{x \in \mathbb{R}^n} \quad & f(x) \\
\text{s.t.} \quad & g_i(x) \le 0, \quad i=1,\dots,m \\
& h_j(x) = 0, \quad j=1,\dots,p
\end{align*}

\subsection*{Skills Required}

You must be able to:

\begin{itemize}
\item Rewrite constraints into standard form $g(x)\le0$
\item Identify decision variables
\item Identify convex vs non-convex structure
\end{itemize}

\subsection*{Common Exam Move}

Turn inequalities like $x_1 \ge 1$ into $1 - x_1 \le 0$.

\newpage

% ============================================================
\section*{[L1-2] Convexity}

\subsection*{Definition}

A function $f$ is convex if:

\[
f(\theta x + (1-\theta)y) \le \theta f(x) + (1-\theta)f(y)
\]

\subsection*{Second Derivative Test}

If twice differentiable:

\[
\nabla^2 f(x) \succeq 0 \quad \Rightarrow \quad f \text{ convex}
\]

\subsection*{Convex Sets}

A set is convex if the line segment between any two points lies in the set.

\subsection*{Exam Skills}

\begin{itemize}
\item Show quadratic forms are convex via positive semidefinite matrices
\item Recognise epigraphs
\item Recognise half-spaces as convex sets
\end{itemize}

\newpage

% ============================================================
\section*{[L1-3] Lagrangian and KKT Conditions}

\subsection*{Lagrangian Construction}

For inequality constraints:

\[
\mathcal{L}(x,\lambda) = f(x) + \sum_i \lambda_i g_i(x)
\]

with $\lambda_i \ge 0$.

\subsection*{KKT Conditions}

\begin{enumerate}
\item Stationarity: $\nabla_x \mathcal{L} = 0$
\item Primal feasibility: $g_i(x)\le0$
\item Dual feasibility: $\lambda_i \ge 0$
\item Complementary slackness: $\lambda_i g_i(x)=0$
\end{enumerate}

\subsection*{Problem-Solving Order}

\begin{enumerate}
\item Write constraints in standard form
\item Build Lagrangian
\item Compute gradients
\item Apply complementary slackness carefully
\item Check feasibility
\end{enumerate}

This is heavily examined.

\newpage

% ============================================================
\section*{[L1-4] Interior-Point and Barrier Concepts}

\subsection*{Barrier Form}

Replace constraint $g(x)\le0$ with:

\[
f(x) - \mu \sum_i \log(-g_i(x))
\]

\subsection*{Key Understanding}

\begin{itemize}
\item $\mu \downarrow 0$ drives solution toward boundary
\item Newton steps solve resulting unconstrained problem
\item Works well for convex problems
\end{itemize}

You must explain why it works conceptually.

\newpage

% ============================================================
\section*{[W1L2-1] Dual Function}

\subsection*{Definition}

\[
q(\lambda) = \inf_x \mathcal{L}(x,\lambda)
\]

Dual problem:

\[
\max_{\lambda \ge 0} q(\lambda)
\]

\subsection*{Weak Duality}

\[
q(\lambda) \le p^\star
\]

Always true.

\subsection*{Strong Duality}

If convex and Slater’s condition holds:

\[
d^\star = p^\star
\]

\subsection*{Exam Skill}

Know how to:

\begin{itemize}
\item Compute infimum analytically
\item Detect when dual becomes $-\infty$
\item Interpret multipliers
\end{itemize}

\newpage

% ============================================================
\section*{[W1L2-2] Slater’s Condition}

If problem is convex and there exists strictly feasible point:

\[
g_i(x) < 0
\]

Then strong duality holds.

You may need to state this precisely.

\newpage

% ============================================================
\section*{[W1L2-3] Dual Interpretation}

Multipliers often represent:

\begin{itemize}
\item Shadow prices
\item Sensitivity of objective to constraint tightening
\end{itemize}

Interpretation questions are common.

\newpage

% ============================================================
\section*{Problem-Solving Checklist}

When given a constrained optimisation problem:

\begin{enumerate}
\item Is the objective convex?
\item Is the feasible set convex?
\item Write constraints in standard form
\item Form Lagrangian
\item Compute gradients
\item Apply KKT conditions
\item Solve systematically
\item Check feasibility
\item If convex, conclude global optimality
\item Optionally construct dual
\end{enumerate}

\newpage

% ============================================================
\section*{Mathematical Skills You Must Be Comfortable With}

\begin{itemize}
\item Multivariable differentiation
\item Matrix algebra
\item Inequality reasoning
\item Logical implication chains
\item Structured equation solving
\end{itemize}

These two lectures are foundational. If you are precise here, later topics (MPC, DP, RL) become structured extensions rather than new mathematics.

\end{document}
